Dynamical systems that undergo a bifurcation of the same type—for example, a Hopf bifurcation—show similar behavior near that transition.
This similarity is captured by normal forms, which are simplified, canonical models to which many different systems can be smoothly transformed.
The center manifold theorem explains why this reduction is possible even for high-dimensional systems: near the bifurcation, the essential dynamics depend only on the variables involved in the stability change, and these variables span the center manifold.
By reducing the system to this low-dimensional manifold, we can describe its behavior using the corresponding normal form.
Here we summarize the main steps to map, close to a Hopf bifurcation, the Wilson-Cowan model
\begin{equation}
    \begin{aligned}
  \tau_x \dot{x}+x &= \gamma_x\,\sigma_x\!\bigl(\kappa_{xx}x - \kappa_{xy}y + P_x\bigr),\\
  \tau_y \dot{y}+y &= \gamma_y\,\sigma_y\!\bigl(\kappa_{yx}x - \kappa_{yy}y\bigr),
\end{aligned}
\end{equation}
to the Stuart-Landau model,
\begin{equation}\label{eq:SLappendix}
\dot{z} \;=\; (\alpha + i\,\omega)\,z \;-\; (\gamma + i\,\beta)\,\lvert z\rvert^{2}z.
\end{equation}
This derivation follows chapter 3.4 from \cite{kuznetsov2023elements}, which contains detailed proofs and explanations.


Let's assume $(x_0,y_0)$ is a fixed point of the Wilson-Cowan model undergoing
a Hopf bifurcation at $P_x=P_c$, letting all the other parameters fixed.
The first step towards the normal form consist of applying
the change of variables $X=x-x_0$ and $Y=y-y_0$, in order to set the
steady state at the origin. Also, we reparametrize $\alpha = P_x-P_c$,
hence
\begin{equation}
    \begin{aligned}
  \tau_x \dot{X}+X+x_0 &= \gamma_x\,\sigma_x\!\bigl(\kappa_{xx}(X+x_0) - \kappa_{xy}(Y+y_0) + P_c + \alpha \bigr),\\
  \tau_y \dot{Y}+Y+y_0 &= \gamma_y\,\sigma_y\!\bigl(\kappa_{yx}(X+x_0) - \kappa_{yy}(Y+y_0)\bigr).
\end{aligned}
\end{equation}
Let 
\begin{equation}
    \begin{aligned}
    I_x &= \kappa_{xx}x_0 - \kappa_{xy}y_0 + P_c + \alpha \quad\text{and}\\
    I_y &= \kappa_{yx}x_0 - \kappa_{yy}y_0. \\
    \end{aligned}
\end{equation}
Then, the Taylor expansion of the velocity fields at the fixed point $(X,Y)=(0,0)$ up to third order reads
\begin{equation}\label{eq:taylor}
    \begin{aligned}
  \frac{\tau_x}{\gamma_x} \dot{X} \approx \;&
  \sigma_x'\!\bigl(I_x \bigr)\left[\kappa_{xx}X - \kappa_{xy}Y\right] + 
  \sigma_x''\!\bigl(I_x \bigr)\left[\frac{\kappa_{xx}^2}{2}X^2 - \kappa_{xx}\kappa_{xy}XY + \frac{\kappa_{xy}^2}{2} Y^2\right] \\
  +\; & \sigma_x'''\!\bigl(I_x \bigr)
  \left[
  \frac{\kappa_{xx}^3}{6}X^3-
  \frac{\kappa_{xx}^2\kappa_{xy}}{2}X^2Y+
  \frac{\kappa_{xx}\kappa{xy}^2}{2}XY^2-
  \frac{\kappa_{xy}^3}{6}Y^3
  \right],\\
  \frac{\tau_y}{\gamma_y} \dot{Y} \approx \;&
  \sigma_y'\!\bigl(I_y \bigr)\left[\kappa_{yx}X - \kappa_{yy}Y\right] + 
  \sigma_y''\!\bigl(I_y \bigr)\left[\frac{\kappa_{yx}^2}{2}X^2 - \kappa_{yx}\kappa_{yy}XY + \frac{\kappa_{yy}^2}{2} Y^2\right] \\
  +\; & \sigma_y'''\!\bigl(I_y \bigr)
  \left[
  \frac{\kappa_{yx}^3}{6}X^3-
  \frac{\kappa_{yx}^2\kappa_{xy}}{2}X^2Y+
  \frac{\kappa_{yx}\kappa{xy}^2}{2}XY^2-
  \frac{\kappa_{yy}^3}{6}Y^3
  \right].
\end{aligned}
\end{equation}
Let's define three generic vectors in $\mathbb{R}^2$: $v_1 = (X_1,Y_2)^T$, $v_2=(X_2,Y_2)^T$, and $v_3=(X_3,Y_3)^T$ 
Now, from the previous Taylor development, we define the Jacobian at the fixed point, $A$, and the multilinear functions $B(v_1,v_2)$ and $C(v_1,v_2,v_3)$ capturing the quadratic and cubic terms of the expansion:
\begin{equation}
A = 
    \begin{pmatrix}
    \sigma_x'\!\bigl(I_x \bigr)\frac{\kappa_{xx}\gamma_x}{\tau_x} & - \sigma_x'\!\bigl(I_x \bigr)\frac{\kappa_{xy}\gamma_y}{\tau_x}\\[5pt]
    \sigma_y'\!\bigl(I_y \bigr)\frac{\kappa_{yx}\gamma_y}{\tau_y}   & - \sigma_x'\!\bigl(I_x \bigr)\frac{\kappa_{yy}\gamma_y}{\tau_y}
    \end{pmatrix},\quad
B(v_1,v_2) =
    \begin{pmatrix}
    B_x(v_1,v_2)\\[5pt]
    B_y(v_1,v_2)
    \end{pmatrix}
    ,\quad
C(v_1,v_2,v_3) =
    \begin{pmatrix}
    C_x(v_1,v_2)\\[5pt]
    C_y(v_1,v_2)
    \end{pmatrix}
\end{equation}
where
\begin{equation}
\begin{aligned}
B_x(v_1,v_2) &= \frac{\gamma_x}{\tau_x}\sigma_x''(I_x)\Big[
\kappa_{xx}^2\,X_1X_2
- \kappa_{xx}\kappa_{xy}\,(X_1Y_2+X_2Y_1)
+ \kappa_{xy}^2\,Y_1Y_2
\Big],\\[6pt]
B_y(v_1,v_2) &=  \frac{\gamma_y}{\tau_y}\sigma_y''(I_y)\Big[
\kappa_{yx}^2\,X_1X_2
- \kappa_{yx}\kappa_{yy}\,(X_1Y_2+X_2Y_1)
+ \kappa_{yy}^2\,Y_1Y_2
\Big].
\end{aligned}
\end{equation}
and
\begin{equation}
\begin{aligned}
C_x(v_1,v_2,v_3) &= 
\frac{\gamma_x}{\tau_x}\sigma_x'''(I_x)\Big[
\kappa_{xx}^3\,X_1X_2X_3
- \kappa_{xx}^2\kappa_{xy}\,(X_1X_2Y_3+X_1X_3Y_2+X_2X_3Y_1)\\
&\qquad
+ \kappa_{xx}\kappa_{xy}^2\,(X_1Y_2Y_3+X_2Y_1Y_3+X_3Y_1Y_2)
- \kappa_{xy}^3\,Y_1Y_2Y_3
\Big],\\[6pt]
C_y(v_1,v_2,v_3) &= 
\frac{\gamma_y}{\tau_y}\sigma_y'''(I_y)\Big[
\kappa_{yx}^3\,X_1X_2X_3
- \kappa_{yx}^2\kappa_{yy}\,(X_1X_2Y_3+X_1X_3Y_2+X_2X_3Y_1)\\
&\qquad
+ \kappa_{yx}\kappa_{yy}^2\,(X_1Y_2Y_3+X_2Y_1Y_3+X_3Y_1Y_2)
- \kappa_{yy}^3\,Y_1Y_2Y_3
\Big].
\end{aligned}
\end{equation}
Notice that Eq. \eqref{eq:taylor} can be expressed as
\begin{equation}
    \dot v = Av + \frac{1}{2}B(v,v) + \frac{1}{6}C(v,v,v)
\end{equation}
where $v=(X,Y)^T$.

With these definitions applying the necessary transformations is relatively straightforward. 
Let
$\lambda_\pm = \mu\pm i\omega$ be the eigenvalues of $A$ near the bifurcation ($\alpha\approx 0$), i.e.,
\begin{equation}
\mu =\frac{\mathrm{tr}(A)}{2}\quad\text{and}\quad\omega =\frac{1}{2} \sqrt{4\det(A) - \mathrm{tr}^2(A)}.
\end{equation}
These are the values of the linear term in the Stuart-Landau model \eqref{eq:SLappendix}.

To obtain the nonlinear parameters $\gamma$ and $\beta$ we need to explicitly compute the left and right eigenvectors associated to $\lambda_+$ and $\overline{\lambda_+}$. Let's name these vectors $p$ and $q$ respectively,
that is $A^T p = \overline{\lambda_+}p$ and $Aq = \lambda_+ q$, with the normalization $\langle p,q \rangle = 1$.
Then, the complex variable for the Hopf normal form is $z = \langle p,v \rangle$,
and its dynamics is given by Eq. \eqref{eq:SLappendix} with 
\begin{equation}
    \gamma + i\beta  = -\frac{g_{20}g_{11}(2\lambda + \bar{\lambda})}{2|\lambda|^2} - \frac{|g_{11}|^2}{\lambda} - \frac{|g_{02}|^2}{2(2\lambda - \bar{\lambda})} - \frac{g_{21}}{2}
\end{equation}
where
\begin{equation}
    g_{20} = \langle p,B(q , q) \rangle, \quad
    g_{11} = \langle p, B(q , \overline{q})\rangle, \quad
    g_{02} = \langle p, B(\overline{q} , \overline{q})\rangle,\quad \text{and}\quad
    g_{21} = \langle p, C(q,q,\overline{q})\rangle.
\end{equation}

